\chapter{坐标变换与几何约束细化}

\section{地固坐标与局部坐标}
为了统一处理不同天体方向，本文采用“地固全局坐标 + 观测方向局部坐标”双坐标系。地固系用于存储节点与地锚点数据，局部系用于描述抛物面轴向几何。设旋转矩阵为
\[
\mathbf R(\alpha,\beta)=\mathbf R_z(\alpha)\mathbf R_y\!\left(\frac{\pi}{2}-\beta\right),
\]
则局部坐标可由 $\mathbf x_{\text{loc}}=\mathbf R\mathbf x_{\text{glob}}$ 得到。

\section{理想抛物面在局部系中的表达}
在局部系下，抛物面可写为标准形式
\[
z=\frac{x^2+y^2}{4f},
\]
再通过逆旋转映射回地固系，实现“求解简化 + 输出一致”。

\section{节点筛选边界}
300 m 口径对应半径 150 m。定义口径筛选条件
\[
\|\mathbf x_i^\perp\|_2\le150,
\]
其中 $\mathbf x_i^\perp$ 为节点在口径平面上的投影坐标。该筛选直接决定优化变量规模和计算时间。

\section{索网边集构造}
由反射面板三角连接关系可还原边集 $E$。边集构造建议去重并排序，避免重复约束导致数值病态。记边集大小为 $|E|$，则边长约束数量与 $|E|$ 同阶。

\section{几何一致性检查}
调节完成后需验证：
\begin{enumerate}
  \item 节点是否仍保持可行拓扑；
  \item 局部法向是否与面板法向方向一致；
  \item 位移是否与促动器径向运动方向相容。
\end{enumerate}
该检查是“数值可行”到“工程可行”的关键桥接步骤。

\section{本章小结}
本章细化了坐标变换与几何约束处理流程，为后续光学计算和鲁棒性分析奠定统一坐标基础。
